\documentclass[11pt]{article}
\usepackage[a4paper, top=18mm, bottom=18mm, left=20mm, right=20mm]{geometry}
\usepackage[T2A]{fontenc}
\usepackage[utf8]{inputenc}
\usepackage[ukrainian]{babel}
\usepackage{graphicx}
\usepackage[export]{adjustbox}
\usepackage{amsmath}
\usepackage{amssymb}
\graphicspath{{/app/Quiz/Scripts/2023/ProgAll/15BinTree/}}
\setlength{\parindent}{0pt}
\setlength{\parskip}{6pt}
\emergencystretch=3em
\pagestyle{empty}
\begin{document}
%begin
{\centering\Large\bfseries Контрольна робота\par}
\medskip
\noindent\rule{\linewidth}{0.5pt}
\smallskip
\noindent\textbf{Варіант \No\,1}\hfill \textit{ПІБ:}\,\underline{\hspace{60mm}}
\vspace{4mm}

%beginitem
1. \textbf{Що буде виведено за виконання фрагменту коду?}

print('R', end=' ')
d=['0','1','2','3','4']
d.insert(1, [1,0])
print(d)

\par\vspace{5mm}
%enditem

%beginitem
2. \textbf{Що буде виведено за виконання фрагменту коду?}

a,b,c=3,4,8
def h(b):
global c    a*=4
    b=5
    c=2
    return a+b+c

a,b,c=6,7,1
print(h(b),a,b,c)
\par\vspace{5mm}
%enditem

%beginitem
3. 

\begin{center}
	\adjustimage{max size={0.8\linewidth}{0.8\paperheight}}
	{../15BinTree/trees_exp/79.png}
\end{center}
Указати послідовність міток вершин дерева, зображеного на рис., що утвориться в результаті обходу дерева в оберненому порядку. Мітки записувати через пробіл.\par Алгоритм починає роботу з кореня (на рис. це верхня вершина).
%answer:79 оберненому
\par\vspace{5mm}
%enditem

%beginitem
4. \textbf{Що буде виведено за виконання фрагменту коду?}
%code
for d in range(-7,-7,2):
    if d>=6:
        break
        print(d,end=' ')
    if d>=6:
        break
else:
    print(d, end=' ')
print(d, end=' ')
%code
\par\vspace{5mm}
%enditem

%beginitem
5. \textbf{Що буде виведено за виконання фрагменту коду?}

%code
print(5.0 >= 5 != False <= 6)
%code
\par\vspace{5mm}
%enditem

%beginitem
6. %goodpath:Scripts/2019/Mod2019_1/Lex/01-good.txt
\textbf{Відмітити все, що є однією правильною лексемою:}
( 1 ) 1__23

( 2 ) 2j

( 3 ) True

( 4 ) 0123
\par\vspace{5mm}
%enditem

%beginitem
7. \textbf{Що буде виведено за виконання фрагменту коду?}
%code
a, b, c = 9, 6, 7
def h(a, b=8, c=6):
    print(a, b, c, end=" ")

h(3, 4, c=5)
print(a, b, c)
%code
\par\vspace{5mm}
%enditem

%beginitem
8. 
Значенням змінної \kw{a} є цілочисловий масив типу $numpy.ndarray$ розмірів $6{\times}6$. На скільки збільшиться сума елементів масиву \kw{a} після виконання
\begin{verbatim}
a.T[0] += 3
\end{verbatim}


Якщо код некоректний, то в якості відповіді зазначити ERROR.


\par\vspace{5mm}
%enditem

%beginitem
9. \textbf{Що буде виведено за виконання фрагменту коду?}

%code
4**1.0-2
%code
\par\vspace{5mm}
%enditem

%beginitem
10. \textbf{Що буде виведено за виконання фрагменту коду?}

%code
cout << (15 % 15 % 15 % 15);
%code
\par\vspace{5mm}
%enditem

%end
\newpage
\end{document}

